\ifdefined\isMainDocument
\else
\documentclass[11pt]{article}
\usepackage{fullpage}
\usepackage{amsmath,amsthm,amsfonts,amssymb,amscd}
\usepackage{bm}
\usepackage{bbm}
\usepackage{mathtools}
\usepackage{mathrsfs}
\usepackage{etoolbox}
\usepackage{nicefrac}
\usepackage{wrapfig}
\usepackage{lineno}
\usepackage{caption}
\usepackage{graphicx}
\graphicspath{}

\usepackage{geometry}
\geometry{top=0.5in, bottom=0.5in, left=0.5in, right=0.5in}
% \geometry{top=1.0in, bottom=1.0in, left=1.0in, right=1.0in}

\usepackage{setspace}
\setstretch{1.5}

% \usepackage{helvet}
% \renewcommand{\familydefault}{\sfdefault}

\usepackage{fancyhdr}
\pagestyle{fancy}
\fancyhf{} % Clear all headers and footers
% \fancyfoot[R]{\thepage} % Set the page number on the right side of the footer
\renewcommand{\headrulewidth}{0pt} % Remove header rule
\setlength{\footskip}{0pt}

\usepackage[backend=bibtex, style=numeric, sorting=none, sortcites=true, maxnames=5]{biblatex}
\DeclareNameAlias{default}{last-first}
\DeclareFieldFormat{pages}{#1}
\DeclareCiteCommand{\cite}
  {\usebibmacro{prenote}}
  {\usebibmacro{citeindex}%
   \printtext[bibhyperref]{\textsuperscript{\textbf{\printfield{labelnumber}}}}}
  {\multicitedelim}
  {\usebibmacro{postnote}}
\renewbibmacro{in:}{}
\renewcommand*{\multicitedelim}{\textsuperscript{,}}
\addbibresource{refs/refs.bib}
\begin{document}
\appendix
\fi

%---------------------------------------------------------------------

\section{Appendix A: Derivation of the Unlinked Baseline}
\label{app:unlinkedapp}
\linenumbers

To derive the baseline reduction in effective population size generated by unlinked heritable fitness variance, we track the cumulative variance in allele frequency change driven by the selected genetic background, independent of continuous spatial linkage.

In a neutral Wright-Fisher population with a constant diploid census size $N$, the variance in allele frequency change per generation due to random genetic drift for a neutral allele at frequency $p$ is the standard gametic sampling variance:

\begin{equation}
    \text{Var}(\Delta p^{\text{drift}}) = \frac{p(1-p)}{2N}
\end{equation}

When variance in reproductive success is heritable ($V_A > 0$), the neutral allele experiences selective interference from the genetic background upon which it resides. In a given generation, a diploid population of $N$ individuals has relative fitnesses $w_i$ with expectation $\mathbb{E}[w_i] = 1$ and variance $\text{Var}(w_i) = V_W$. The frequency of the neutral allele in the subsequent generation is approximately:

\begin{equation}
    p' \approx p + \frac{1}{2N}\sum_{i=1}^{N}(w_i - 1)z_i
\end{equation}
where $z_i \in \{0, \frac{1}{2}, 1\}$ denotes the allele copy number in individual $i$ under Hardy-Weinberg proportions. Because the neutral locus is unlinked from all selected sites, relative fitness $w_i$ and the neutral genotype $z_i$ are uncorrelated in expectation ($\mathbb{E}[(w_i-1)z_i] = 0$). The variance in allele frequency change driven purely by this single generation of heritable fitness variation is:

\begin{equation}
    \text{Var}(\Delta p^{\text{sel}}) =\frac{1}{4N^2}\sum_{i=1}^{N}\text{Var}((w_i - 1)z_i) \approx V_A \frac{p(1-p)}{2N}
\end{equation}

This formulation (Robertson 1961) relies on a large-$N$ approximation, assumes Hardy-Weinberg equilibrium ($\text{Var}(z_i) = p(1-p)/2$), and isolates the additive genetic component $V_A$ from the total fitness variance, as contributions from dominance and epistasis do not transmit stably across generations to drive multi-generational allele frequency changes.

Because the neutral allele and the selected background are unlinked (segregating on non-homologous chromosomes), they assort independently during meiosis ($r = 0.5$). Consequently, the statistical association, or linkage disequilibrium, between the neutral allele and its initial selected background halves each generation. The directional change in the neutral allele frequency driven by this specific selective cohort decays geometrically:
\begin{equation}
    \Delta p_t^{\text{sel}} = \left(\frac{1}{2}\right)^t \Delta p^{\text{sel}}
\end{equation}

Integrating this effect over time yields the total cumulative change in the neutral allele frequency resulting from a single generation of heritable variance:

\begin{equation}
    \Delta p_{\text{total}}^{\text{sel}} = \sum_{t=0}^{\infty} \left(\frac{1}{2}\right)^t \Delta p^{\text{sel}} = 2 \Delta p^{\text{sel}}
\end{equation}

Assuming $V_A$ remains stationary under mutation-selection-drift equilibrium, a new influx of selective variance of magnitude $V_A p(1-p)/(2N)$ is generated each generation. The long-run footprint of each generation's selective variation resolves to $2\Delta p^{\text{sel}}$.

The total selective variance in allele frequency change is the variance of this cumulative sum. Because the selective association persists across multiple generations, the variance of the cumulative effect scales with the square of the cumulative sum:
\begin{equation}
    \text{Var}(\Delta p_{\text{total}}^{\text{sel}}) = \text{Var}(2 \Delta p^{\text{sel}}) = 4 V_A \frac{p(1-p)}{2N}
\end{equation}
This establishes that the multi-generational persistence of associations across independently assorting chromosomes amplifies the effective variance four-fold.

Because the neutral allele is physically unlinked from the causal selected loci, the variance generated by selective interference and the variance generated by neutral demographic sampling arise from independent biological processes. Their contributions to the total variance in allele frequency are therefore additive. Equating the standard definition of effective gametic sampling variance to the sum of these independent variances yields:
\begin{equation}
    \frac{p(1-p)}{2N_e} = \text{Var}(\Delta p^{\text{drift}}) + \text{Var}(\Delta p_{\text{total}}^{\text{sel}})
\end{equation}

Substituting the expanded terms:
\begin{equation}
    \frac{p(1-p)}{2N_e} = \frac{p(1-p)}{2N} + 4 V_A \frac{p(1-p)}{2N}
\end{equation}

Dividing by the common neutral variance term $p(1-p)$ and multiplying by 2 yields the unlinked expectation for the effective population size:
\begin{equation}
    \frac{1}{N_e} = \frac{1}{N} (1 + 4V_A) \Rightarrow N_e = \frac{N}{1 + 4V_A}
\end{equation}
This baseline establishes the minimum selective interference generated by heritable fitness variance, irrespective of spatial chromosomal architecture. The extension to separate sexes and structured mating systems is derived in Appendix B, where the additive variance $V_A$ is scaled by the mating system parameter $\kappa(\alpha)$, yielding $N_e/N = 1/(1+4\kappa V_A)$.

\ifdefined\isMainDocument
\else
    \end{document}
\fi